%==========================================================================================
\section{\texttt{xrf\_matrix}: matrix effects in layered samples}\label{sec:matrix}
%==========================================================================================

\subsection{Overview and physics}

The fluorescence intensity of an element is proportional to its concentration only in an
infinitely thin sample. In a real sample the incident beam is attenuated on its way in, the
fluorescence on its way out (\emph{primary absorption}), and the fluorescence of one element
can excite another (\emph{secondary fluorescence}). The module computes both for a stack of
homogeneous layers---cap, film, substrate---of arbitrary thickness, the last one possibly
infinitely thick, under polychromatic excitation. It follows the fundamental-parameter
approach of \citet{Sherman1955} and \citet{Shiraiwa1966}, extended to multilayers
\citep{Mantler1986,deBoer1990}; see also the handbook of \citet{VanGrieken2002}.

\paragraph{Notation.} Layers are numbered from the surface, $k = 0,1,\dots$; layer $k$ has
density $\rho_k$, thickness $t_k$ and mass thickness $M_k = \rho_k t_k$ (infinite for a bulk
last layer). The beam meets the surface at the angle $\psi_\mathrm{in}$, the detector sees it
at the take-off angle $\psi_\mathrm{out}$, both measured from the surface. $\mu_k(E)$ is the
total mass attenuation coefficient of layer $k$, $Q_{il}(E)$ the partial fluorescence
cross-section of line $l$ of element $i$ (xraylib \fn{CS_FluorLine_Kissel}, cm$^2$/g:
photo-ionisation of the subshell, Coster--Kronig transfers, fluorescence yield and radiative
rate), $w_i^k$ the mass fraction of $i$ in layer $k$, and $\dd\Phi/\dd E$ the incident
density per keV. Write
\begin{equation}
  a_k(E) = \frac{\mu_k(E)}{\sin\psi_\mathrm{in}}, \qquad
  b_k = \frac{\mu_k(E_l)}{\sin\psi_\mathrm{out}}, \qquad
  F(c;M) = \int_0^M \e^{-cm}\dd m = \frac{1-\e^{-cM}}{c}, \qquad
  G(x,y;M) = \int_0^M \e^{-xm}\e^{-y(M-m)}\dd m .
\end{equation}

\paragraph{Primary fluorescence.} A photon of energy $E$ reaching mass depth $m$ of layer
$k$ has crossed $m/\sin\psi_\mathrm{in}$ of it; the probability that it produces line $l$ of
element $i$ in the slab $\dd m$ is $w_i^k Q_{il}(E)\,\dd m/\sin\psi_\mathrm{in}$, and the
fluorescence photon reaching the detector has crossed $m/\sin\psi_\mathrm{out}$. Integrating
over the depth and over the spectrum, and leaving out the solid angle and the detection
efficiency (common to all lines and applied by \fn{xrf_fit.Detector.efficiency}), the
sensitivity per unit mass fraction is
\begin{equation}
  P_{il}^k = T_\mathrm{out}^k \int \frac{\dd\Phi}{\dd E}\,\frac{Q_{il}(E)}{\sin\psi_\mathrm{in}}\,
  T_\mathrm{in}^k(E)\,F\!\left(a_k(E)+b_k;\,M_k\right)\dd E,
  \label{eq:primary}
\end{equation}
with the overlayer transmissions $T_\mathrm{in}^k(E) = \exp(-\sum_{c<k}a_c M_c)$ and
$T_\mathrm{out}^k = \exp(-\sum_{c<k}b_c M_c)$. In the thin limit $F \to M_k$ (no
absorption); for a bulk layer $F = 1/(a_k+b_k)$, Sherman's equation.

\paragraph{Secondary fluorescence.} Line $j$ (energy $E_j$) of an element of layer $a$,
produced by the beam at depth $m'$, is emitted isotropically. From a plane source the
fraction absorbed at a distance $x$ (in mass thickness) per unit thickness is
$\frac12\int_0^1 \mu\,\e^{-\mu x/u}\,\dd u/u$ with $u$ the direction cosine; replacing $\mu$
by the partial cross-section $w_i^k Q_{il}(E_j)$ gives the rate of secondary emission of line
$l$. With $\kappa_c = \mu_c(E_j)/u$ the depth integrals are closed-form for every $u$:
\begin{align}
  S_{il}^k &= T_\mathrm{out}^k \sum_{a,j} \tfrac12\,Q_{il}(E_j)\,w_j^a \int
  \frac{\dd\Phi}{\dd E}\,\frac{Q_j(E)}{\sin\psi_\mathrm{in}}\,T_\mathrm{in}^a(E)
  \int_0^1 K_{ak}(E,u)\,\frac{\dd u}{u}\,\dd E, \label{eq:secondary}\\
  a = k:&\quad K = D(a_k, b_k, \kappa_k; M_k) = \int_0^{M_k}\!\!\int_0^{M_k}
  \e^{-a_k m' - b_k m - \kappa_k|m-m'|}\,\dd m\,\dd m', \notag\\
  a < k:&\quad K = G(a_a,\kappa_a;M_a)\;\e^{-\sum_{a<c<k}\kappa_c M_c}\;F(b_k+\kappa_k;M_k),
  \notag\\
  a > k:&\quad K = F(a_a+\kappa_a;M_a)\;\e^{-\sum_{k<c<a}\kappa_c M_c}\;G(b_k,\kappa_k;M_k).
  \notag
\end{align}
The same-layer double integral has the closed form
\begin{equation}
  D(a,b,\kappa;M) = \frac{F(a+\kappa)-F(a+b)}{b-\kappa} + \frac{F(b+\kappa)-F(a+b)}{a-\kappa},
  \label{eq:D}
\end{equation}
(both terms written with $F(\cdot) = F(\cdot\,;M)$), obtained by splitting the domain at
$m=m'$; when $b\to\kappa$ or $a\to\kappa$ the corresponding quotient becomes $-F'(a+\kappa)$
or $-F'(b+\kappa)$. For a single bulk layer Eq.~\eqref{eq:secondary} reduces to the classical
$\left[\ln(1+a/\mu_j)/a + \ln(1+b/\mu_j)/b\right]/[2(a+b)]$ \citep{Shiraiwa1966}; for a thin
layer it gives the logarithmic thin-film enhancement. The integral over $u$ is computed on a
logarithmic grid ($\int_0^1 f\,\dd u/u = \int f\,\dd(\ln u)$), which resolves the grazing paths
that dominate in thin films.

The measured intensity of line $l$ of element $i$ is then
\begin{equation}
  I_{il} \propto \epsilon(E_l)\sum_k w_i^k\left(P_{il}^k + S_{il}^k\right).
\end{equation}

\begin{hyp}
\begin{itemize}[nosep,leftmargin=1em]
\item Homogeneous layers with flat, parallel, sharp interfaces and infinite lateral extent;
  no roughness, interdiffusion or lateral inhomogeneity.
\item A single incidence and take-off angle (point-like beam footprint seen from a distant
  detector).
\item No refraction and no total external reflection: at incidence angles comparable to the
  critical angle (a few mrad) the penetration of the beam is not the geometrical one
  assumed here.
\item Secondary fluorescence from the lines of the stack's own elements only (contaminants
  do not excite), lines above 1\,keV, and lines contributing at least \fn{exciter_min_rel}
  of their element's strongest line. No tertiary fluorescence and no excitation by
  scattered radiation or by photo- or Auger electrons.
\item The secondary term is smooth in $E$ and is evaluated on a decimated energy grid.
\end{itemize}
\end{hyp}

\pycode{xrf_matrix}{_header}

\subsection{Cached cross-sections}

\api[\texttt{\_grid\_key}, \texttt{\_cs\_total}, \texttt{\_cs\_fluor}, \texttt{\_cs\_fluor\_scalar}]{\texttt{\_grid\_key(e)}, \texttt{\_cs\_total(z, e\_kev)}, \texttt{\_cs\_fluor(z, macro, e\_kev)}, \texttt{\_cs\_fluor\_scalar(z, macro, e\_kev)}}
xraylib is called one energy at a time, so the cross-sections on the spectrum's energy grid
are computed once and cached, keyed by the element, the line and a hash of the grid.
\fn{_cs_fluor} returns zero below the edge (where xraylib raises \fn{ValueError}) and clips
negative values. \fn{_IUPAC_LINES} holds the IUPAC line macros only (see
Section~\ref{sec:lines} for why the Siegbahn names are excluded).
\pycode{xrf_matrix}{_grid_key}
\pycode{xrf_matrix}{_cs_total}
\pycode{xrf_matrix}{_cs_fluor}
\pycode{xrf_matrix}{_cs_fluor_scalar}

\subsection{Geometry and sample}

\api[\texttt{Geometry}]{\texttt{Geometry(incidence\_deg=40.0, takeoff\_deg=40.0)}}
Incidence and take-off angles in degrees from the surface, and their sines (properties
\fn{sin_in}, \fn{sin_out}). Both default to 40\textdegree; set them to the geometry of the
measurement, since the path lengths scale as $1/\sin\psi$.
\pycode{xrf_matrix}{Geometry}
\pycode{xrf_matrix}{Geometry.sin_in}
\pycode{xrf_matrix}{Geometry.sin_out}

\api[\texttt{Layer}]{\texttt{Layer(composition, density\_g\_cm3, thickness\_um=None, name="")}}
A homogeneous layer.
\begin{args}
\item[composition] a formula (\fn{"In0.5Ga0.5As"}) or \texttt{\{symbol: atoms\}}; only atomic
  proportions matter.
\item[density\_g\_cm3] density.
\item[thickness\_um] thickness; \fn{None} or \fn{inf} for a bulk layer.
\item[name] label.
\end{args}
\blk{Attributes} \fn{formula}, \fn{mass_fractions} (\texttt{\{Z: w\}}, from xraylib's
\fn{CompoundParser}), \fn{atoms} (\texttt{\{symbol: n\}}), \fn{density_g_cm3},
\fn{thickness_um}. \fn{mass_thickness} is $\rho t$ in g/cm$^2$ (infinite for bulk);
\fn{mu(e_kev)} is $\mu(E) = \sum_i w_i(\mu/\rho)_i(E)$, cm$^2$/g, with the cache for grids of
more than 8 points; \fn{with_composition} and \fn{with_thickness} return modified copies.
\pycode{xrf_matrix}{Layer}
\pycode{xrf_matrix}{Layer.__init__}
\pycode{xrf_matrix}{Layer.mass_thickness}
\pycode{xrf_matrix}{Layer.mu}
\pycode{xrf_matrix}{Layer.with_composition}
\pycode{xrf_matrix}{Layer.with_thickness}
\pycode{xrf_matrix}{Layer.__repr__}

\api[\texttt{Stack}]{\texttt{Stack(layers)}}
Ordered layers, surface first; only the last may be bulk (\fn{ValueError} otherwise). Supports
\fn{len}, indexing, iteration, and \fn{replace(index, layer)}, which returns a new stack.
\pycode{xrf_matrix}{Stack}
\pycode{xrf_matrix}{Stack.__init__}
\pycode{xrf_matrix}{Stack.__len__}
\pycode{xrf_matrix}{Stack.__getitem__}
\pycode{xrf_matrix}{Stack.__iter__}
\pycode{xrf_matrix}{Stack.replace}
\pycode{xrf_matrix}{Stack.__repr__}

\subsection{Depth integrals}

\api[\texttt{\_F}, \texttt{\_dF}]{\texttt{\_F(c, M)}, \texttt{\_dF(c, M)}}
$F(c;M) = (1-\e^{-cM})/c$ computed as \fn{-expm1(-cM)/c}, which keeps full precision for
$cM \ll 1$ (thin layers), and $1/c$ for $M=\infty$; \fn{_dF} is its derivative with respect
to $c$, $[cM\e^{-cM}-(1-\e^{-cM})]/c^2$ ($-1/c^2$ for bulk), used in the limits of
Eq.~\eqref{eq:D}.
\pycode{xrf_matrix}{_F}
\pycode{xrf_matrix}{_dF}

\api[\texttt{\_G}]{\texttt{\_G(x, y, M)}}
$G(x,y;M) = \e^{-\min(x,y)M}\,(1-\e^{-|x-y|M})/|x-y|$, a form that never evaluates the
exponential of a large positive number, with the limit $M\,\e^{-xM}$ for $x=y$.
\pycode{xrf_matrix}{_G}

\api[\texttt{\_D}]{\texttt{\_D(a, b, k, M, rtol=1e-4)}}
Eq.~\eqref{eq:D}, broadcasting over arrays (energies $\times$ direction cosines), switching to
the derivative form when $|b-\kappa|$ or $|a-\kappa|$ is within \fn{rtol} of their sum.
\pycode{xrf_matrix}{_D}

\subsection{\texttt{MatrixModel}}

\api{\texttt{\_Exciter}}
An exciting line: layer index, element, line name, energy (keV), mass fraction, its
cross-section $Q_j(E)$ on the secondary-fluorescence grid, and $\mu_c(E_j)$ in every layer.
\pycode{xrf_matrix}{_Exciter}

\api[\texttt{MatrixModel}]{\texttt{MatrixModel(sample, geometry, energy\_kev, flux\_01bw, secondary=True, target=None, sf\_step\_kev=0.04, n\_u=160, exciter\_min\_rel=1e-3)}}
Primary and secondary sensitivities of every line in a layer stack, for one incident
spectrum and geometry.
\begin{args}
\item[sample] a \fn{Layer} (one-layer stack) or a \fn{Stack}.
\item[geometry] a \fn{Geometry}.
\item[energy\_kev, flux\_01bw] incident spectrum per 0.1\,\% bandwidth, as returned by
  \fn{Beamline.spectrum} or \fn{monochromatic_spectrum}; converted to per keV internally. A
  uniform grid with $\sim$10\,eV steps is appropriate for a white beam.
\item[secondary] include secondary fluorescence.
\item[target] index of the layer being quantified (required for a multilayer); per-layer
  quantities default to it, and elements absent from the stack are treated as if in it.
\item[sf\_step\_kev] step of the decimated grid for the secondary term (default 40\,eV).
\item[n\_u] points of the logarithmic direction-cosine grid ($u$ from $10^{-7}$ to 1).
\item[exciter\_min\_rel] relative threshold for exciting lines.
\end{args}
The constructor precomputes $\mu_k(E)$ of every layer on the grid, the decimated grids and
the incident transmissions $T_\mathrm{in}^k(E)$.
\pycode{xrf_matrix}{MatrixModel}
\pycode{xrf_matrix}{MatrixModel.__init__}

\api[\texttt{layer}, \texttt{with\_layer}, \texttt{with\_stack}]{\texttt{layer} (property), \texttt{with\_layer(layer, index=None)}, \texttt{with\_stack(stack, target=None)}}
The target layer, and copies of the model with one layer or the whole stack replaced (same
beam, geometry, numerics and target). Used by the self-consistent quantification.
\pycode{xrf_matrix}{MatrixModel.layer}
\pycode{xrf_matrix}{MatrixModel.with_layer}
\pycode{xrf_matrix}{MatrixModel.with_stack}

\api[\texttt{exciters}]{\texttt{exciters} (property)}
The lines of the stack's elements that can drive secondary fluorescence, computed once: for
each element of each layer, every line above 1\,keV, ranked by its primary production at the
top of its layer, $\int\Phi\,Q_j\,T_\mathrm{in}^a\,\dd E/\sin\psi_\mathrm{in}$; lines below
\fn{exciter_min_rel} of the strongest line of that element are dropped.
\pycode{xrf_matrix}{MatrixModel.exciters}

\api[\texttt{\_k}, \texttt{\_b}, \texttt{\_tout}]{\texttt{\_k(layer)}, \texttt{\_b(line\_energy\_kev)}, \texttt{\_tout(b, k)}}
Helpers: the layer index (target by default); the exit attenuation $b_c = \mu_c(E_l)/
\sin\psi_\mathrm{out}$ of every layer; and $T_\mathrm{out}^k$.
\pycode{xrf_matrix}{MatrixModel._k}
\pycode{xrf_matrix}{MatrixModel._b}
\pycode{xrf_matrix}{MatrixModel._tout}

\api[\texttt{primary}]{\texttt{primary(z, macro, line\_energy\_kev, layer=None)}}
Eq.~\eqref{eq:primary}, integrated with the trapezoidal rule on the full grid. Returns 0 if
the line cannot be excited by the spectrum.
\pycode{xrf_matrix}{MatrixModel.primary}

\api[\texttt{secondary}]{\texttt{secondary(z, macro, line\_energy\_kev, layer=None, by\_exciter=False)}}
Eq.~\eqref{eq:secondary}: for each exciter able to excite the line
($Q_{il}(E_j) > 0$), the kernel is computed on the (energy, $u$) grid, integrated over
$\ln u$, then over energy. With \fn{by_exciter=True} it also returns the contribution of each
exciter, keyed \fn{"<layer>:<element> <line>"}---useful to see which element drives the
enhancement.
\pycode{xrf_matrix}{MatrixModel.secondary}

\api{\texttt{sensitivity}, \texttt{emission}, \texttt{enhancement}, \texttt{absorption\_factor}}
\fn{sensitivity} is $P^k + S^k$ per unit mass fraction in layer $k$. \fn{emission} is the
intensity from the whole stack, $\sum_k w_i^k(P^k [+S^k])$; it is what sets relative
intensities \emph{within} a line group in \fn{xrf_fit} (elements absent from the stack are
evaluated per unit mass fraction in the target). \fn{enhancement} is $S/P$.
\fn{absorption_factor} is $P^k$ divided by the absorption-free value $M_k\int\Phi
Q\,\dd E/\sin\psi_\mathrm{in}$ (thin film, no overlayer): 1 means no absorption; NaN for a
bulk layer.
\pycode{xrf_matrix}{MatrixModel.sensitivity}
\pycode{xrf_matrix}{MatrixModel.emission}
\pycode{xrf_matrix}{MatrixModel.enhancement}
\pycode{xrf_matrix}{MatrixModel.absorption_factor}

\subsection{Quantification}

\api[\texttt{\_line}]{\texttt{\_line(group)}}
\fn{(Z, macro, energy in keV)} of a group's reference line; merged lines (\fn{"KL3+1"}) are
refused, since their cross-section is not defined.
\pycode{xrf_matrix}{_line}

\api[\texttt{quantify\_layer}]{\texttt{quantify\_layer(areas, groups, matrix, detector, reference, analytes, fixed\_atoms=None, max\_iter=50, tol=1e-6)}}
Self-consistent composition of the target layer from fitted reference-line areas.
\begin{args}
\item[areas] \texttt{\{label: reference-line area\}}, e.g. \fn{FitResult.areas}.
\item[groups] \texttt{\{label: LineGroup\}}.
\item[matrix] a \fn{MatrixModel}; its target layer gives the starting composition.
\item[detector] an \fn{xrf_fit.Detector} (detection efficiency).
\item[reference] label of the reference group (\fn{"As K"}).
\item[analytes] labels of the groups to quantify.
\item[fixed\_atoms] \texttt{\{symbol: atoms per reference atom\}} for unmeasured elements of the
  layer (e.g. oxygen).
\item[max\_iter, tol] iteration limit and convergence threshold on the atom ratios.
\end{args}
\blk{Method} For each group $g$ the measured area is modelled as
\begin{equation}
  A_g \propto w_g^t s_g^t + C_g, \qquad s_g^k = (P^k+S^k)\,\epsilon, \qquad
  C_g = \sum_{k\neq t} w_g^k s_g^k ,
\end{equation}
where $t$ is the target layer and $C_g$ the known contribution of the other layers (caps,
substrate). Relative to the reference group $r$,
\begin{equation}
  w_i^t = \frac{(A_i/A_r)\,(w_r^t s_r^t + C_r) - C_i}{s_i^t} .
\end{equation}
The atom ratios $n_i/n_r = (w_i^t/W_i)/(w_r^t/W_r)$ ($W$: atomic weights) rebuild the target
composition, the sensitivities are recomputed, and the loop stops when no ratio changes by
more than \fn{tol}. Density and thicknesses are held fixed. Elements of the target that are
neither analytes, reference nor fixed are dropped, because the absolute mass fraction of the
reference matters once other layers contribute. A negative $w_i$ (the other layers explain
more than the measured intensity) is clipped to 0 with a warning. For a one-layer sample
$C=0$ and the result is $n_i/n_r = (A_i/A_r)(s_r W_r)/(s_i W_i)$, the formula of
\fn{xrf_fit.number_ratio} with matrix-corrected sensitivities.
\blk{Returns} a dictionary: \fn{ratios} (atoms per reference atom), the converged
\fn{layer} and \fn{matrix}, \fn{k} (the factors $(n_i/n_r)/(A_i/A_r)$), \fn{enhancement} and
\fn{absorption} of the reference lines, \fn{other_layers} (fraction of each modelled
intensity coming from the other layers), \fn{terms} ($s^t$, $C$), \fn{iterations},
\fn{history}, \fn{reference}.
\pycode{xrf_matrix}{quantify_layer}

\api[\texttt{atom\_ratio}]{\texttt{atom\_ratio(quant, label, area, area\_ref, err=None, err\_ref=None)}}
Applies the converged sensitivities of \fn{quantify_layer} to other areas---typically every
pixel of a map---with the matrix held fixed:
$n_i/n_r = [(R\,I_r - C_i)/s_i/W_i]/(w_r/W_r)$, $R = A_i/A_r$, $I_r = w_r s_r + C_r$. Errors of
the two areas are propagated as uncorrelated. Returns \fn{(ratio, error)}.
\pycode{xrf_matrix}{atom_ratio}
\pycode{xrf_matrix}{_label_symbol}
\pycode{xrf_matrix}{_ref_symbol}
